% IEEE conference-format draft. Compile on Overleaf (upload paper/ and results/figures/) or: pdflatex main && bibtex main && pdflatex main x2
% Convention: \todo{...} marks text that still needs real numbers/results. Nothing red should remain at submission.
\documentclass[conference]{IEEEtran}
\usepackage{cite,amsmath,graphicx,booktabs,xcolor,url,hyperref}
\graphicspath{{../results/figures/}}
\newcommand{\todo}[1]{\textcolor{red}{[TODO: #1]}}
% include a figure only if the file exists, so the draft compiles before results are generated
\newcommand{\safefig}[3]{\IfFileExists{../results/figures/#1.png}{\includegraphics[width=#2]{#1}}{\fbox{\todo{run pipeline: #1.png}}}}

\begin{document}
\title{Drift-Aware Retraining for Credit-Card Fraud Detection:\\A Chronological Comparison of Retraining Policies on IEEE-CIS}
\author{\IEEEauthorblockN{Aditya Katkar, Member~2, Member~3, Member~4}
\IEEEauthorblockA{Department of Computer Engineering\\\todo{Institute name}\\Guide: Dr.~Kapil Vhatkar}}
\maketitle

\begin{abstract}
Fraud-detection models are usually evaluated on a single static train/test split, which hides how they behave once deployed
and fraud tactics change. We replay the IEEE-CIS Fraud Detection data chronologically in weekly windows, quantify how a model that
is never retrained degrades, evaluate ADWIN and DDM as drift detectors, and compare three retraining policies---no retraining,
fixed-interval and drift-triggered---on predictive performance and operational cost.
\todo{one sentence per headline result: decay size, static-vs-deployment gap, best policy, cost saving}
\end{abstract}
\begin{IEEEkeywords}
fraud detection, concept drift, model retraining, MLOps, ADWIN, DDM, LightGBM
\end{IEEEkeywords}

\section{Introduction}
Card-not-present fraud evolves continuously: attackers adapt to countermeasures and legitimate customer behaviour shifts with seasons
and products. A model trained once therefore loses accuracy after deployment, a phenomenon known as concept
drift~\cite{gama2014survey}. Most published fraud-detection studies nevertheless report results on a static random split, and rarely
examine \emph{when} a deployed model should be retrained. We address this with a chronological simulation of deployment and study:
\begin{itemize}
\item \textbf{RQ1} How much does a model trained once degrade under chronological deployment, and how optimistic is a static split?
\item \textbf{RQ2} Can ADWIN and DDM detect the degradation, and with what delay and false-alarm behaviour?
\item \textbf{RQ3} Which retraining policy gives the best performance-versus-cost trade-off?
\item \textbf{RQ4} Which features drive predictions, and does their importance shift with drift?
\end{itemize}
\textbf{Contributions.} (i) A reproducible chronological-deployment benchmark on IEEE-CIS with an explicit label-delay option;
(ii) a quantification of the gap between static-split and deployment-time performance; (iii) a cost-aware comparison of
no, fixed-interval and drift-triggered retraining; (iv) an open-source MLOps pipeline coupling a drift detector to retraining.
\todo{final contribution list once End-Sem results are in}

\section{Related Work}
\textbf{Concept drift.} Gama et al.~\cite{gama2014survey} survey drift adaptation. DDM~\cite{gama2004ddm} monitors the online error rate and
signals drift when it rises significantly above its historical minimum; ADWIN~\cite{bifet2007adwin} keeps an adaptive window and
shrinks it when two sub-windows differ significantly.
\textbf{Fraud detection under drift.} Dal Pozzolo et al.~\cite{dalpozzolo2018} show that supervised information in fraud detection arrives
delayed, so feedback and delayed samples must be treated differently; Adebayo et al.~\cite{adebayo2023} review machine-learning and
drift techniques for card fraud. \todo{2--3 more papers from the literature review (Member 4); compare on: dataset, drift handling, delay, cost}
\textbf{Gap.} \todo{state precisely what prior work does not do: cost-aware comparison of retraining policies with label delay on IEEE-CIS}

\section{Data and Problem Formulation}
We use the IEEE-CIS Fraud Detection data~\cite{ieeecis}: transaction and identity tables joined on \texttt{TransactionID}; target
\texttt{isFraud}. Only the training files are labelled, so all experiments use them. \texttt{TransactionDT} is a time offset in seconds;
we use it only for ordering and windowing. Table~\ref{tab:data} summarises the data; 12 columns with more than 90\% missing values were dropped.

\begin{table}[t]\centering\caption{Dataset statistics.}\label{tab:data}
\begin{tabular}{ll}\toprule Statistic & Value\\\midrule
Transactions & 590{,}540\\ Fraudulent transactions & 20{,}663 (3.50\%)\\ Time span (days) & 182\\
Transactions with identity record & 24.4\%\\ Missing values (raw, merged) & 45.0\%\\
Columns: raw merged / after cleaning & 435 / 423\\ Columns used (after correlation pruning) & 304\\ \bottomrule\end{tabular}
\end{table}

\textbf{Deployment simulation.} Transactions are sorted by time and split into $W$ windows of 7 days. A model is trained on the first
$w_0=8$ windows. For each later window $w$ it scores the window using only data available before $w$; a retraining policy
then decides whether to retrain. With label delay $\delta$, a transaction becomes a training example only $\delta$ days after it occurs.
The decision threshold is chosen on the last training window (F1-optimal) and never on deployment data.

\section{Methodology}
\subsection{Preprocessing}
Columns with more than 90\% missing values and constant columns are dropped; numeric columns are downcast; categorical columns are
stored as categories with one global mapping (no target information). Missing values are left to the tree learners.
Highly correlated numeric features ($|r|>0.95$) are pruned using the initial training window only. Class imbalance is handled by
class weighting; Section~\ref{sec:imb} compares it with undersampling and SMOTE~\cite{chawla2002smote}.
\subsection{Model}
LightGBM~\cite{ke2017lightgbm} is the main learner; XGBoost~\cite{chen2016xgboost} is used for comparison. \todo{final hyper-parameters (Member 4 tuning)}
\subsection{Drift detection}
ADWIN and DDM from River~\cite{montiel2021river} run over two binary streams of the never-retrained model: errors on all transactions, and misses on
fraud transactions only (closer to production, where confirmed fraud arrives via chargebacks). Covariate shift is measured by the population stability index (PSI), the
Kolmogorov--Smirnov statistic, and the AUC of an adversarial classifier separating the reference period from a later week.
\subsection{Retraining policies}
\emph{No retraining}; \emph{fixed interval} (retrain every $k$ windows); \emph{drift-triggered} (retrain when a detector fires).
\todo{End-Sem: detector choice, cool-down, sliding vs expanding training window}
\subsection{Cost model}
\todo{End-Sem: $\text{Cost}= c_{FN}\sum \text{missed fraud amount} + c_{FP} N_{FP} + c_{R} N_{\text{retrains}}$, with justified constants and a sensitivity analysis}
\subsection{Explainability}
\todo{End-Sem: SHAP global importance per window; compare before/after drift}

\section{Experimental Setup}
Metrics: ROC-AUC, PR-AUC (primary, because of the severe class imbalance), recall at 1\% false-positive rate, and threshold metrics (precision, recall, F1).
\todo{hardware, software versions, random seed, run time}

\section{Results}
\subsection{RQ1: degradation without retraining}
Trained on the first 8 weeks and never updated, the LightGBM model's weekly PR-AUC falls from 0.546 (mean of the first three
deployment weeks) to 0.471 (last three), a 13.7\% relative drop (Spearman $\rho=-0.49$, $p=0.039$; slope $-0.006$ per week).
ROC-AUC falls from 0.892 to 0.858 (3.8\%, $\rho=-0.60$, $p=0.008$). Recall at 1\% false-positive rate (6.6\% drop) and F1 (7.1\%) also
decline but their trends are not statistically significant ($p=0.32$ and $0.22$), so we treat the degradation as real but moderate.
The evaluation protocol matters more than the drift itself: the same learner scores PR-AUC 0.721 (ROC-AUC 0.954) on a random 80/20 split,
0.526 on a temporal 80/20 split, and 0.478 when deployed once and scored on all later weeks (0.452 on the last quarter of weeks),
i.e.\ a static split overstates PR-AUC by about one third (Fig.~\ref{fig:prot}). LightGBM outperformed XGBoost under the same protocol
(mean weekly PR-AUC 0.470 vs 0.428).
\begin{figure}[t]\centering\safefig{decay_curve}{\columnwidth}{decay}\caption{Performance of the never-retrained model by deployment week.}\end{figure}
\begin{figure}[t]\centering\safefig{evaluation_protocols}{0.95\columnwidth}{prot}\caption{The same learner under different evaluation protocols.}\label{fig:prot}\end{figure}
\subsection{Class-imbalance handling}\label{sec:imb}
Trained on the initial window and scored on every later week, plain training reached a mean PR-AUC of 0.523, class weighting 0.498,
random undersampling (10:1) 0.517 and undersampling followed by SMOTE 0.514. Imbalance handling therefore did not help on this data;
we keep class weighting in the simulator for consistency with the literature but report this negative result. \todo{re-run the simulator without class weights as a sensitivity check}
\subsection{RQ2: drift detection}
\begin{figure}[t]\centering\safefig{drift_detection}{\columnwidth}{drift}\caption{Detector alarms against the measured degradation onset.}\end{figure}
\begin{figure}[t]\centering\safefig{psi_heatmap}{\columnwidth}{psi}\caption{PSI of the most important features over deployment time.}\end{figure}
Run on the error stream of the never-retrained model, all four detector configurations (ADWIN and DDM, on all transactions or on
fraud misses only) fire in the first two deployment weeks and repeatedly afterwards, nine to ten weeks \emph{before} the measured
degradation onset (week 18), so as raw monitors they behave mostly as false alarms. Covariate drift, by contrast, is large and monotone:
the AUC of an adversarial classifier separating the training period from week $w$ rises from 0.77 to 0.97, three of the twenty most
important features exceed PSI 0.25 by the last week, and the most shifted feature is the browser version \texttt{id\_31} (PSI 1.69).
Drift in the inputs is therefore much clearer than drift in performance, which only declines moderately (Section~RQ1).
\subsection{RQ3: retraining policies}
Table~\ref{tab:pol} compares the policies over the 18 deployment weeks. Every retraining policy improves on never retraining
(paired Wilcoxon on weekly PR-AUC, $p\le0.008$ for all but the 14-day label-delay variant). Performance rises with retraining frequency:
fixed 2-weekly retraining gives the best mean PR-AUC (0.546, +16.2\%), followed by 4-weekly (0.530, +12.8\%) and 8-weekly (0.511, +8.8\%).
Drift-triggered retraining did \emph{not} beat fixed intervals at a comparable number of retrainings: ADWIN on the error stream retrained four
times and reached 0.534, matching fixed 4-weekly retraining (0.530), and DDM on fraud misses reached 0.523 with three retrainings versus 0.511 for
two fixed retrainings. We therefore do not find support for the hypothesis that drift triggers retrain more efficiently than a well-chosen fixed interval on this data.
Under the default cost assumptions all retraining policies are cheaper than never retraining (by 3.7\% to 16.2\%); the saving is driven by fraud
losses, not by retraining effort. The break-even retraining cost, i.e.\ the price per retraining at which a policy stops paying for itself, is
at least \$17k for every policy against an assumed \$500, so retraining pays off under any plausible price. Label delay erodes the benefit:
for 4-weekly retraining the mean PR-AUC falls from 0.530 (no delay) to 0.516 (7 days) and 0.496 (14 days, $p=0.060$ against never retraining).
Results come from a single run on one dataset, so differences of a few percent between policies should not be over-interpreted.

\begin{table}[t]\centering\caption{Retraining policies over 18 deployment weeks (cost at \$2 per review, \$500 per retraining).}\label{tab:pol}
\small\begin{tabular}{lrrrr}\toprule Policy & Retrains & PR-AUC & Gain & Cost saving\\\midrule
Never retrain & 0 & 0.470 & -- & --\\ Fixed 2 weeks & 8 & 0.546 & +16.2\% & 16.2\%\\ Fixed 4 weeks & 4 & 0.530 & +12.8\% & 16.0\%\\
Fixed 8 weeks & 2 & 0.511 & +8.8\% & 14.9\%\\ ADWIN (errors) & 4 & 0.534 & +13.5\% & 15.0\%\\ ADWIN (misses) & 1 & 0.497 & +5.7\% & 3.7\%\\
DDM (errors) & 4 & 0.520 & +10.5\% & 10.1\%\\ DDM (misses) & 3 & 0.523 & +11.2\% & 10.8\%\\ PR-AUC drop rule & 2 & 0.518 & +10.3\% & 12.2\%\\
\bottomrule\end{tabular}\end{table}
\begin{figure}[t]\centering\safefig{policy_tradeoff}{\columnwidth}{tradeoff}\caption{Performance versus retraining effort, and cost breakdown per policy.}\end{figure}
\subsection{RQ4: explainability}
Exact TreeSHAP values of the deployed model show a \emph{stable} reliance on the same signals: the importance ranking of the first and
last three deployment weeks is almost identical (Spearman $\rho=0.995$, 9 of the top-10 features shared). The notable exceptions are the
identity strings \texttt{id\_31} and \texttt{id\_30} (browser and operating-system versions), which lose 12 and 10 rank positions. They are also
the most shifted features by PSI (\texttt{id\_31}: 1.69), which suggests that new browser/OS versions unseen in training erode
a useful signal rather than the model's logic changing. \todo{interpretation of risers (D11, id\_02, M3) once confirmed}

\section{Discussion}
\todo{Implications for practitioners; covariate vs concept drift; effect of label delay}
\section{Threats to Validity}
Single dataset; \texttt{TransactionDT} has no calendar dates, so seasonality cannot be identified; anonymised features limit interpretability;
cost constants are assumptions; detector hyper-parameters may favour one method. \todo{add after experiments}
\section{Conclusion}
\todo{3--4 sentences answering RQ1--RQ4}

\section*{Acknowledgment}
We thank Dr.~Kapil Vhatkar for guidance.

\bibliographystyle{IEEEtran}
\bibliography{references}
\end{document}
